\section{Discussion and Limitations}\label{discussion}
The main unresolved issue is the gap between predicted observation value and the joint task outcome. Softer structural evidence restores the influence of a finite category prior, and native branches show that this influence can change actual surface acquisition. Yet an object-focused gain can displace a larger whole-scene gain or delay unknown coverage. The complete left-layout comparison demonstrates that this tradeoff can persist beyond the first changed target. A better probability loss is therefore useful model evidence, not a substitute for coverage and reconstruction evaluation.

Geometry remains conditional and imperfect. The particle model integrates depth and placement while retaining estimated orientation, width and parts. Ideal instance labels do not remove these geometric errors. The earlier fixed-state audits identify unassigned fragments receiving whole templates, unreliable outer-plane fits, unobserved depth and imperfect layer placement; their failures remain in the supplement. Independent procedural measurements have exact poses and noiseless depth, and the two exploration layouts were inspected during development. Neither independent object tests nor multiple policy runs convert these related floorplan permutations into broad environment generalization.

The class-temperature adjustment and split-conformal interval correction are standard calibration tools applied to distinct targets. Positive temperature cannot correct the sign of neutral structural evidence, while wide residual intervals can improve coverage without improving usefulness. The interval diagnostic's mean width exceeds 2\,m and its proper interval score worsens. Nominal interval coverage, class calibration and navigation utility must consequently remain separate claims. The matched G control also shares ideal instance association; it should not be interpreted as a semantics-free perception baseline.

Computational and physical limits remain substantial. Policies use a logical action budget and exact virtual poses; three concurrent runs with numerical-library thread limits of one do not establish real-time robot control. Native measured reconstruction is evaluated against a fixed source reference, without predicted geometry insertion or post-hoc alignment. The physical records lack such a reference and include localization/control failures. Deployment with estimated poses, natural detections and robust execution requires a separate matched evaluation after hardware recovery.

\section{Conclusion}\label{conclusion}
We have formulated and tested joint unknown-space coverage and surface observation with a category-conditioned structural reward. The implementation preserves the separation between acquired maps, hypothetical geometry and execution outcomes. Its score decomposition makes semantic inactivity diagnosable: changing a class probability matters only when structural observation opportunities and route margins permit a different target.

Independent structural calibration reduces probability loss, and complete simulated runs establish full-coverage feasibility under the declared action and return contract. Category information can change native target selection, but the complete comparison remains adverse in one permutation and tied in the other. The evidence therefore supports an active-planning system and a concrete semantic decision mechanism, while leaving a general coverage or reconstruction advantage unproven. Future improvements should be judged by complete unknown-exploration outcomes alongside object and whole-scene surface acquisition, rather than by stronger probabilities or selected favorable branches.

\section{Reproducibility and Data}\label{reproducibility-and-data}
The primary policy traces to the frozen v10\_3\_1 observed-state scoring/runtime source and its declared joint-exploration adapter. The common coverage and measured-quality terms, original candidates, heading-state local planner and execution contract are retained in the compared arms. Separate frozen modules specify conditional geometry and structural-temperature changes. Historical interface names OV-SDF, STGHP, RPN-UQ and IGCR do not imply that a trained reachability network or complete original supplied-graph controller was active in these virtual runs.

The accompanying sources contain object-level splits, generated sensor manifests, sanitized actor inputs, frozen method/configuration hashes, all candidate scores, acquired action records, native meshes and independent evaluation receipts. Public summary tables identify reused controls and branches. No test-set score, hidden source mesh or future observation is supplied to the actor. Python file-access audit hooks provide provenance checks for trusted code; they are not an operating-system security sandbox. The supplement preserves earlier controlled results, null/negative exploration comparisons, rejected geometry variants and physical failures. Author identities and submission details remain unfilled.
